% ═══════════════════════════════════════════════════════════════════
% ЧАСТЬ I (продолжение). ИСТОРИЯ И МЕТОДОЛОГИЯ ПРОГРАММЫ
% ═══════════════════════════════════════════════════════════════════

\section{История программы: от поправок к замкнутой системе}\label{sec:history}

Программа развивалась в несколько этапов, каждый из которых добавлял
новый слой строгости. Понимание этой эволюции помогает читателю
оценить, почему окончательная редакция выглядит именно так: все
методологические решения продиктованы конкретными контрольными
событиями, зафиксированными протоколами.

\subsection{Этап I: три поправки}

Исходным ядром послужил набор из трёх алгебраических поправок
(раздел~\ref{sec:corrections}): спиновая фаза $\delta=\pi/N$,
торможение $\gamma=\delta^4/k$ и эффективная фаза
$\delta_{\mathrm{eff}}=\delta^5/k$ с объединённой формулой
\eqref{eq:deltaCh}. Уже на этом этапе действовало правило
целочисленности: каждая величина обязана восходить к целочисленной
характеристике геометрии. Правило проверялось на единственной тогда
ступени --- кривой Клейна, где тройка $(7,22,\lambda_1)$ бралась из
группы $\Gamma(2,3,7)$ и второй коприсоединённой K3.

\subsection{Этап II: калибровочные стенды}

Второй этап добавил два стенда-калибровки: тор и K3. Тор дал полный
явный вывод поправок из первых принципов (раздел~\ref{sec:torus}) и
обнаружил важный факт: эквивариантность $\mu_4$ для симметричной
кодировки \emph{точная}, а для кодировки $(7,22)$ --- нарушенная.
Этот контраст показал, что поправки наследуют геометрию носителя,
и был сформулирован принцип естественности: поправка обязана
коммутировать с группой симметрии. Принцип получил доказательство в
сертификатах~C и~D. Стенд K3 (раздел~\ref{sec:k3}) дал первую
крупную точную выкладку: ранг $20$, сигнатура $(1,19)$, определитель
$64$ --- всё над $\QQ$, с полным контролем.

\subsection{Этап III: сертификаты}

Третий этап ввёл институт сертификатов: каждый нетривиальный переход
конвейера обязан иметь независимую проверку либо точной арифметикой,
либо прямым интегрированием. Появились сертификатор циклов (A) и
сертификат~B, закрывший ядро размерности $29$ периодами. За ними
последовали C (естественность через $\mu_4$), D (универсальная
$\mu_4$-теорема), E (завершимость потока с точным временем),
F (рационализация с априорной границей знаменателей) и G (индексы
решётки Ходжа, слой Коула--Селберга). Институт сертификатов оказался
наиболее продуктивным элементом программы: он превращает каждое
«вроде бы верно» в «проверено с точностью $10^{-36}$» или
«доказано тождеством».

\subsection{Этап IV: циклотомический стенд и замкнутая система}

Четвёртый этап --- стенд $N=15/30$ (часть~V). Здесь программа
достигла полной зрелости: изложение замкнуто в системе лемм и теорем,
не опирающейся ни на PSLQ, ни на L-значения, ни на численные проверки
доказуемого. Все фазы выведены явно (лемма~\ref{lem:closedform}),
соотношения Римана доказаны тождественно
(теорема~\ref{th:riemann}), нормировка приведена к дробным долям
$\pi$ (теорема~\ref{th:normalization}), индексы решётки Ходжа
вычислены точно (теорема~\ref{th:indexcomp},
сертификат~H). Сертификат~H довершил арку: полная $\Ga$-нормировка
Гросса объёма дала $\mathrm{vol}_h=1125=\sqrt{\disc}$ --- целое
число, красивое в своей арифметике $3^4\cdot5^6$.

\subsection{Errata: два контрольных события}

История программы включает два честно зафиксированных исправления,
которые сами стали результатами. Первое --- errata E8 (раздел
\ref{sec:e8}): на роде $3$ чётных спиновых структур $36$, нечётных $28$ —
а не обратная пара, как было напечатано в ранней версии; полное перечисление
$64$ структур по инварианту Арфа закрыло вопрос навсегда и добавило
в программу автоматический переборный модуль. Второе --- кейс E6
на торе: корректный вывод поправок даёт $\Delta=0.7990$, а не
ранее напечатанное $0.5964$; исправление получено той же цепочкой
первых принципов и внесено во все таблицы. Оба события
проиллюстрировали силу протокольного подхода: расхождение между
напечатанным и вычисленным обнаруживается автоматически при
воспроизведении.

\section{Методология: пять правил программы}\label{sec:rules}

Сформулируем пять правил, которыми программа руководствуется с
этапа~III. Правила просты, но именно они определяют архитектуру
монографии и лаборатории.

\begin{keybox}{Пять правил программы}
\begin{enumerate}[leftmargin=2em, itemsep=0.2em]
\item \emph{Правило тройки.} Геометрия поставляет только
  $(n,B,\lambda_0)$. Всё остальное выводится.
\item \emph{Правило замкнутой формы.} Каждое проверяемое число
  либо выводится в замкнутой алгебраической форме, либо
  сертифицируется прямым независимым измерением.
\item \emph{Правило независимости сертификата.} Сертификат не может
  использовать тот же алгоритм, что и основное вычисление; аналитика
  сверяется интегрированием без $\Ga$-функций, комбинаторика ---
  точной арифметикой.
\item \emph{Правило воспроизводимости.} Каждый результат
  воспроизводится одной командой из чистого репозитория; случайность
  и сиды запрещены.
\item \emph{Правило фиксации.} Расхождения фиксируются как errata с
  полным анализом; исправление входит в основной текст, а не в
  сноски.
\end{enumerate}
\end{keybox}

Правило~1 отделяет геометрию от алгебры и делает фреймворк
переносимым: новая геометрия требует только тройки. Правило~2
устраняет «подгонку»: число либо выводится, либо измерено независимо.
Правило~3 защищает от замкнутого круга: сертификат, использующий те
же функции, что и проверяемое вычисление, ничего не сертифицирует.
Правило~4 делает монографию и лабораторию единым организмом: каждое
число текста имеет источник в \texttt{laboratory.py}. Правило~5
превращает ошибки в результаты: errata E8 и кейс E6 --- полноценные
разделы с автоматической проверкой.

\subsection{Почему вычислительные сертификаты дополняют доказательства}

Может возникнуть вопрос: зачем сертификаты, если теоремы доказаны?
Ответ состоит в разделении труда между аналитическим и вычислительным
каналом. Теорема фиксирует тождество в общем виде; сертификат
гарантирует, что конкретная реализация --- код, библиотеки, точность
--- действительно вычисляет то самое тождество, а не близкое к нему.
Ошибка в знаке, сдвиг индекса, потеря ветви комплексного логарифма ---
вот класс ошибок, которые не видны в формулировке теоремы, но
немедленно проявляются в независимом численном сверке. На стенде
$N=15/30$ сертификат~B ловит именно это: замкнутая форма
$\frac1N\zeta^{ra+sb}\Omega_{a,b}$ и независимое интегрирование
согласуются до $10^{-36}$, значит, реализация корректна, а фазы
произведены в нужном порядке.

Симметричная ситуация возникает с целочисленными структурами: тип
поляризации $(1,1,5,5,15,15,15,15)$ получен алгоритмом Смита, но
его арифметика (произведение инвариантных множителей $= \disc$,
совпадение с $3^4\cdot5^6$) проверена отдельной целочисленной
строкой. Двойная регистрация --- алгоритм плюс арифметический
инвариант --- правило программы для всех ключевых таблиц.
